% !TeX root = ../se200_referential_regimes.tex

% =============================================================================
\section{Background: Neutral Substrates}
\label{sec:background}
% =============================================================================

This paper builds on \emph{Neutral Substrates}~\citep{case2026neutral}.
That work studies a shared representational base intended to support
reference under persistent interpretive disagreement without adopting
contested causal or normative interpretations as unqualified foundational
commitments.
The present paper investigates proposed identity conditions
and their behavior when representations are transformed.
It does not
assume that the abstract referential interfaces have been given adequate
identity criteria for every intended kind of referent.

The three definitions below retain their SE-100 labels and wording.
They establish the terminology imported for the present analysis;
they do not independently justify a concrete referential regime.

% -----------------------------------------------------------------------------
\subsection{Inherited Definitions}
\label{subsec:ns-inherited-definitions}
% -----------------------------------------------------------------------------

\begin{definition}[Substrate]
  \label{se100.def.Substrate}
  A \emph{substrate} $\Substrate$ is a shared representational base providing
  stable reference for entities, occurrences, and institutional artifacts
  across a class of admissible interpretive frameworks $\Frameworks$.
  Stable reference consists of individuation, co-reference, and persistence.
\end{definition}

A substrate is the foundational layer, not the complete record system.
It can support later causal, normative, evidentiary, or institutional
interpretations without adopting those interpretations as its own
object-level commitments.
The distinction between a referent and
interpretations attributed to it is central; preserving a shared label
or descriptive record does not itself establish that the intended
referent has been fixed.

\begin{definition}[Referential Regime]
  \label{se100.def.ReferentialRegime}
  A \emph{referential regime} is the triple of
  \begin{enumerate}
    \item individuation conditions,
    \item co-reference conditions, and
    \item persistence conditions
  \end{enumerate}
  by which a substrate fixes and tracks entities, occurrences, and institutional
  artifacts.
\end{definition}

Individuation concerns what counts as one whole and sameness across
presentations; co-reference concerns whether references denote the same
referent; persistence concerns strict identity across settings and must
be distinguished from succession, derivation, or counterpart relations.
The abstract triple supplies interfaces for those questions, not their
independently warranted answers.
A relation defined by closure of
proposed preserving steps does not, without additional grounding,
establish identity in an intended domain.

\begin{definition}[Referential Commitments]
  \label{se100.def.ReferentialCommitments}
  Let $\SubstrateRef$ denote the \emph{referential commitments} of $\Substrate$.
  These are the substrate-layer commitments fixed by its referential regime,
  including identifiers;
  the typing of entities, occurrences, and institutional artifacts;
  timestamps;
  provenance;
  and referential relations among them.
\end{definition}

Identifiers, types, timestamps, provenance, and relations may supply
materials for referential fixing.
Their mere presence does not justify
which object-language propositions are fixed by a proposed referential
regime.
That connection requires independently grounded interpretation.

\emph{Neutral Substrates} distinguishes a permitted attribution from
the proposition attributed: accepting the attribution does not require
accepting its content.
It also defines interpretive non-commitment,
extension stability, semantic neutrality, and neutrality by design.

Its central result is conditional sufficiency.
Under Referential Common
Ground, restricting foundational commitments to referential commitments
and permitted attribution propositions suffices for semantic neutrality
relative to every admissible framework.
Permitted attributions are
determined by the referential commitments.

The result does not establish that Referential Common Ground holds in
a given realization, that the adopted referential criteria are adequate,
that neutrality has been guaranteed by a permitted design-time basis,
or that the foundational restriction is necessary for every neutral
architecture.
In particular, the formal restriction alone does not
show that no object-level causal or normative proposition is contained
in the purported referential commitments.
Those boundaries also apply
to the present paper.

% -----------------------------------------------------------------------------
\subsection{Contribution}
\label{subsec:contribution}
% -----------------------------------------------------------------------------

This paper selects a finite reference inventory, candidate identity
bases, and transformation families through which to investigate
identity-sensitive distinctions.
It studies classifications of
re-expression, annotation, refinement, aggregation or decomposition,
re-contextualization, reassignment, substitution, branching,
provenance extension, and state evolution.
The resulting classification
profiles support limited formal comparison; their empirical or semantic
adequacy remains an independent question.

\begin{definition}[Transformation-Invariant Identity]
  \label{se200.def.TransformationInvariantIdentity}
  Relative to a specified referent kind, proposed identity basis, and
  admissible transformation instances, a presentation is treated as
  \emph{transformation-invariant} along steps classified as preserving
  the declared basis.
  A transformation classification expresses a
  proposed identity judgment relative to that basis; it does not,
  without independent justification, establish strict identity of
  the represented referent across those steps.
\end{definition}

Section~\ref{sec:formal_setting} formalizes candidate profiles and a
relation generated from non-breaking steps;
Section~\ref{sec:algebra} compares selected basis-sensitive classifications.
The counting
question and the need for independent scientific warrant are examined
in Sections~\ref{sec:distinctions} and~\ref{sec:limits}.
